% --- Archon physics formalization source begin ---
% archon:physics
% archon:covers PhyXMiniProblems/problem_phyx_mini_0947.lean
% archon:source-report reports/phyx_mini/problem_phyx_mini_0947.source.json

\chapter{Physics problem phyx\_mini\_0947}
\label{ch:PhyXMiniProblems_problem_phyx_mini_0947}

\paragraph{Problem source.}
\#\# Physical scenario

The rectangular loop is pivoted about the y-axis and carries a current of \textbackslash{}( 15.0\textbackslash{},\textbackslash{}text\{A\} \textbackslash{}) in the direction indicated.

\#\# Question

If the loop is in a uniform magnetic field with magnitude \textbackslash{}( 0.48\textbackslash{},\textbackslash{}text\{T\} \textbackslash{}) in the \textbackslash{}( +x \textbackslash{})-direction, find the magnitude and direction of the torque required to hold the loop in the position shown.

\#\# Answer choices

- A: A: 0.042N·m

- B: B: 0.030N·m

- C: C: 0.026N·m

- D: D: 0.034N·m

\#\# Figure caption (auxiliary; use the image as primary evidence)

The image depicts a three-dimensional diagram with a rectangular loop oriented in a specific position within an XYZ coordinate system. Here are the details:

1. **Objects and Structure:**
   - A rectangular loop is shown, outlined with a yellow border.
   - The rectangle is tilted in space, with its left side aligned along the y-axis and its base along the x-z plane.

2. **Dimensions and Angles:**
   - The left side of the rectangle is labeled as 8.00 cm.
   - The bottom side of the rectangle is labeled as 6.00 cm.
   - The bottom side makes a 30.0° angle with the negative x-axis.

3. **Current Flow:**
   - Arrows indicate the direction of current flow around the loop.
   - A specific segment of the loop on the left indicates a current of 15.0 A moving upward.

4. **Coordinate Axes:**
   - The y-axis is vertical and depicted going upwards.
   - The x-axis is horizontal and pointing rightwards.
   - The z-axis is oriented into the page from a three-dimensional perspective.

The diagram illustrates the orientation and dimensions of a rectangular current-carrying loop within a coordinate system.

\#\# Dataset metadata

Magnetism; Physical Model Grounding Reasoning, Spatial Relation Reasoning

\paragraph{Recorded answer/context.}
B: B: 0.030N·m

\paragraph{Figure/image path.}
/root/proposal\_for\_physic/science-mango/phyx\_mini\_run/phyx\_data/test\_image/947.png

\paragraph{Formalization target.}
create a compiling Lean file with sorry bodies at `PhyXMiniProblems/problem\_phyx\_mini\_0947.lean`. The Lean declarations must preserve the physical quantities, dimensions or dimensional roles, figure labels, governing-law hypotheses, and final relation expressed by this problem.
Use Mathlib/Physlib names found through LeanExplore where available. If a physics API is missing, introduce faithful local abstractions rather than scalar placeholder aliases.

\begin{theorem}[Physics formalization target]
\label{thm:physics:phyx_mini_0947:target}
\lean{PhyXMiniProblems.ProblemPhyXMini0947.problem_phyx_mini_0947}
\uses{def:physics:phyx-mini-0947:phyxminiproblems-problemphyxmini0947-torqueinnewtonmeters, def:physics:phyx-mini-0947:phyxminiproblems-problemphyxmini0947-torquemagnitudeinnewtonmeters, def:physics:phyx-mini-0947:phyxminiproblems-problemphyxmini0947-yhat, def:physics:phyx-mini-0947:phyxminiproblems-problemphyxmini0947-rectangularcurrentloopsetup, def:physics:phyx-mini-0947:phyxminiproblems-problemphyxmini0947-matchesproblemstatement, def:physics:phyx-mini-0947:phyxminiproblems-problemphyxmini0947-matchesprimaryfigure, def:physics:phyx-mini-0947:phyxminiproblems-problemphyxmini0947-matchesrectangularloopgeometry, def:physics:phyx-mini-0947:phyxminiproblems-problemphyxmini0947-modelsuniformmagneticfield, def:physics:phyx-mini-0947:phyxminiproblems-problemphyxmini0947-satisfiescurrentloopmagneticmomentlaw, def:physics:phyx-mini-0947:phyxminiproblems-problemphyxmini0947-satisfiesmagnetictorquelaw, def:physics:phyx-mini-0947:phyxminiproblems-problemphyxmini0947-satisfiesstaticholdingcondition, def:physics:phyx-mini-0947:phyxminiproblems-problemphyxmini0947-pointsalongpositiveaxis, def:physics:phyx-mini-0947:phyxminiproblems-problemphyxmini0947-recordedanswerchoice, def:physics:phyx-mini-0947:phyxminiproblems-problemphyxmini0947-agreesatmillinewtonmeterprecision}
The assigned autoformalize agent should translate this physics problem into Lean declarations in the covered file, with theorem and lemma proof bodies written as `by sorry`.
\end{theorem}
\begin{proof}
This is an autoformalization task, not a proof task. Produce statements that can later be proved by the physics prover without weakening the physical model.
\end{proof}

% --- Archon declaration topology begin ---
\section{Declaration topology}
\begin{definition}[Spatial Vector]
  \label{def:physics:phyx-mini-0947:phyxminiproblems-problemphyxmini0947-spatialvector}
  \lean{PhyXMiniProblems.ProblemPhyXMini0947.SpatialVector}
  Three-dimensional Euclidean vectors used for coherent-unit readouts
\end{definition}
\begin{proof}
  This is the declared model definition of Spatial Vector.
\end{proof}
\begin{definition}[electric Current Dimension]
  \label{def:physics:phyx-mini-0947:phyxminiproblems-problemphyxmini0947-electriccurrentdimension}
  \lean{PhyXMiniProblems.ProblemPhyXMini0947.electricCurrentDimension}
  Electric current has physical dimension charge per time
\end{definition}
\begin{proof}
  This is the declared model definition of electric Current Dimension.
\end{proof}
\begin{definition}[magnetic Flux Density Dimension]
  \label{def:physics:phyx-mini-0947:phyxminiproblems-problemphyxmini0947-magneticfluxdensitydimension}
  \lean{PhyXMiniProblems.ProblemPhyXMini0947.magneticFluxDensityDimension}
  Magnetic flux density (tesla) has dimension M T⁻¹ C⁻¹
\end{definition}
\begin{proof}
  This is the declared model definition of magnetic Flux Density Dimension.
\end{proof}
\begin{definition}[magnetic Moment Dimension]
  \label{def:physics:phyx-mini-0947:phyxminiproblems-problemphyxmini0947-magneticmomentdimension}
  \lean{PhyXMiniProblems.ProblemPhyXMini0947.magneticMomentDimension}
  \uses{def:physics:phyx-mini-0947:phyxminiproblems-problemphyxmini0947-electriccurrentdimension}
  A current-loop magnetic moment has dimension current times area
\end{definition}
\begin{proof}
  This is the declared model definition of magnetic Moment Dimension.
\end{proof}
\begin{definition}[torque Dimension]
  \label{def:physics:phyx-mini-0947:phyxminiproblems-problemphyxmini0947-torquedimension}
  \lean{PhyXMiniProblems.ProblemPhyXMini0947.torqueDimension}
  Torque has the energy dimension M L² T⁻²
\end{definition}
\begin{proof}
  This is the declared model definition of torque Dimension.
\end{proof}
\begin{definition}[Length Magnitude]
  \label{def:physics:phyx-mini-0947:phyxminiproblems-problemphyxmini0947-lengthmagnitude}
  \lean{PhyXMiniProblems.ProblemPhyXMini0947.LengthMagnitude}
  A nonnegative, unit-independent physical length
\end{definition}
\begin{proof}
  This is the declared model definition of Length Magnitude.
\end{proof}
\begin{definition}[Electric Current Magnitude]
  \label{def:physics:phyx-mini-0947:phyxminiproblems-problemphyxmini0947-electriccurrentmagnitude}
  \lean{PhyXMiniProblems.ProblemPhyXMini0947.ElectricCurrentMagnitude}
  \uses{def:physics:phyx-mini-0947:phyxminiproblems-problemphyxmini0947-electriccurrentdimension}
  A nonnegative, unit-independent electric-current magnitude
\end{definition}
\begin{proof}
  This is the declared model definition of Electric Current Magnitude.
\end{proof}
\begin{definition}[Magnetic Flux Density Magnitude]
  \label{def:physics:phyx-mini-0947:phyxminiproblems-problemphyxmini0947-magneticfluxdensitymagnitude}
  \lean{PhyXMiniProblems.ProblemPhyXMini0947.MagneticFluxDensityMagnitude}
  \uses{def:physics:phyx-mini-0947:phyxminiproblems-problemphyxmini0947-magneticfluxdensitydimension}
  A nonnegative, unit-independent magnetic-flux-density magnitude
\end{definition}
\begin{proof}
  This is the declared model definition of Magnetic Flux Density Magnitude.
\end{proof}
\begin{definition}[Magnetic Flux Density Vector]
  \label{def:physics:phyx-mini-0947:phyxminiproblems-problemphyxmini0947-magneticfluxdensityvector}
  \lean{PhyXMiniProblems.ProblemPhyXMini0947.MagneticFluxDensityVector}
  \uses{def:physics:phyx-mini-0947:phyxminiproblems-problemphyxmini0947-spatialvector, def:physics:phyx-mini-0947:phyxminiproblems-problemphyxmini0947-magneticfluxdensitydimension}
  A unit-independent magnetic-flux-density vector
\end{definition}
\begin{proof}
  This is the declared model definition of Magnetic Flux Density Vector.
\end{proof}
\begin{definition}[Magnetic Moment Vector]
  \label{def:physics:phyx-mini-0947:phyxminiproblems-problemphyxmini0947-magneticmomentvector}
  \lean{PhyXMiniProblems.ProblemPhyXMini0947.MagneticMomentVector}
  \uses{def:physics:phyx-mini-0947:phyxminiproblems-problemphyxmini0947-spatialvector, def:physics:phyx-mini-0947:phyxminiproblems-problemphyxmini0947-magneticmomentdimension}
  A unit-independent magnetic-dipole-moment vector
\end{definition}
\begin{proof}
  This is the declared model definition of Magnetic Moment Vector.
\end{proof}
\begin{definition}[Torque Vector]
  \label{def:physics:phyx-mini-0947:phyxminiproblems-problemphyxmini0947-torquevector}
  \lean{PhyXMiniProblems.ProblemPhyXMini0947.TorqueVector}
  \uses{def:physics:phyx-mini-0947:phyxminiproblems-problemphyxmini0947-spatialvector, def:physics:phyx-mini-0947:phyxminiproblems-problemphyxmini0947-torquedimension}
  A unit-independent torque vector
\end{definition}
\begin{proof}
  This is the declared model definition of Torque Vector.
\end{proof}
\begin{definition}[nonnegative SIReadout]
  \label{def:physics:phyx-mini-0947:phyxminiproblems-problemphyxmini0947-nonnegativesireadout}
  \lean{PhyXMiniProblems.ProblemPhyXMini0947.nonnegativeSIReadout}
  Read a nonnegative dimensionful scalar in coherent SI units
\end{definition}
\begin{proof}
  This is the declared model definition of nonnegative SIReadout.
\end{proof}
\begin{definition}[vector SIReadout]
  \label{def:physics:phyx-mini-0947:phyxminiproblems-problemphyxmini0947-vectorsireadout}
  \lean{PhyXMiniProblems.ProblemPhyXMini0947.vectorSIReadout}
  \uses{def:physics:phyx-mini-0947:phyxminiproblems-problemphyxmini0947-spatialvector}
  Read a dimensionful spatial vector in coherent SI units
\end{definition}
\begin{proof}
  This is the declared model definition of vector SIReadout.
\end{proof}
\begin{definition}[length In Meters]
  \label{def:physics:phyx-mini-0947:phyxminiproblems-problemphyxmini0947-lengthinmeters}
  \lean{PhyXMiniProblems.ProblemPhyXMini0947.lengthInMeters}
  \uses{def:physics:phyx-mini-0947:phyxminiproblems-problemphyxmini0947-lengthmagnitude, def:physics:phyx-mini-0947:phyxminiproblems-problemphyxmini0947-nonnegativesireadout}
  Read a length in metres
\end{definition}
\begin{proof}
  This is the declared model definition of length In Meters.
\end{proof}
\begin{definition}[area In Square Meters]
  \label{def:physics:phyx-mini-0947:phyxminiproblems-problemphyxmini0947-areainsquaremeters}
  \lean{PhyXMiniProblems.ProblemPhyXMini0947.areaInSquareMeters}
  \uses{def:physics:phyx-mini-0947:phyxminiproblems-problemphyxmini0947-nonnegativesireadout}
  Read an area in square metres
\end{definition}
\begin{proof}
  This is the declared model definition of area In Square Meters.
\end{proof}
\begin{definition}[current In Amperes]
  \label{def:physics:phyx-mini-0947:phyxminiproblems-problemphyxmini0947-currentinamperes}
  \lean{PhyXMiniProblems.ProblemPhyXMini0947.currentInAmperes}
  \uses{def:physics:phyx-mini-0947:phyxminiproblems-problemphyxmini0947-electriccurrentmagnitude, def:physics:phyx-mini-0947:phyxminiproblems-problemphyxmini0947-nonnegativesireadout}
  Read an electric-current magnitude in amperes
\end{definition}
\begin{proof}
  This is the declared model definition of current In Amperes.
\end{proof}
\begin{definition}[magnetic Flux Density In Teslas]
  \label{def:physics:phyx-mini-0947:phyxminiproblems-problemphyxmini0947-magneticfluxdensityinteslas}
  \lean{PhyXMiniProblems.ProblemPhyXMini0947.magneticFluxDensityInTeslas}
  \uses{def:physics:phyx-mini-0947:phyxminiproblems-problemphyxmini0947-magneticfluxdensitymagnitude, def:physics:phyx-mini-0947:phyxminiproblems-problemphyxmini0947-nonnegativesireadout}
  Read a magnetic-flux-density magnitude in teslas
\end{definition}
\begin{proof}
  This is the declared model definition of magnetic Flux Density In Teslas.
\end{proof}
\begin{definition}[magnetic Flux Density Vector In Teslas]
  \label{def:physics:phyx-mini-0947:phyxminiproblems-problemphyxmini0947-magneticfluxdensityvectorinteslas}
  \lean{PhyXMiniProblems.ProblemPhyXMini0947.magneticFluxDensityVectorInTeslas}
  \uses{def:physics:phyx-mini-0947:phyxminiproblems-problemphyxmini0947-spatialvector, def:physics:phyx-mini-0947:phyxminiproblems-problemphyxmini0947-magneticfluxdensityvector, def:physics:phyx-mini-0947:phyxminiproblems-problemphyxmini0947-vectorsireadout}
  Read a magnetic-flux-density vector in tesla coordinates
\end{definition}
\begin{proof}
  This is the declared model definition of magnetic Flux Density Vector In Teslas.
\end{proof}
\begin{definition}[magnetic Moment In Ampere Square Meters]
  \label{def:physics:phyx-mini-0947:phyxminiproblems-problemphyxmini0947-magneticmomentinamperesquaremeters}
  \lean{PhyXMiniProblems.ProblemPhyXMini0947.magneticMomentInAmpereSquareMeters}
  \uses{def:physics:phyx-mini-0947:phyxminiproblems-problemphyxmini0947-spatialvector, def:physics:phyx-mini-0947:phyxminiproblems-problemphyxmini0947-magneticmomentvector, def:physics:phyx-mini-0947:phyxminiproblems-problemphyxmini0947-vectorsireadout}
  Read a magnetic moment in ampere-square-metre coordinates
\end{definition}
\begin{proof}
  This is the declared model definition of magnetic Moment In Ampere Square Meters.
\end{proof}
\begin{definition}[torque In Newton Meters]
  \label{def:physics:phyx-mini-0947:phyxminiproblems-problemphyxmini0947-torqueinnewtonmeters}
  \lean{PhyXMiniProblems.ProblemPhyXMini0947.torqueInNewtonMeters}
  \uses{def:physics:phyx-mini-0947:phyxminiproblems-problemphyxmini0947-spatialvector, def:physics:phyx-mini-0947:phyxminiproblems-problemphyxmini0947-torquevector, def:physics:phyx-mini-0947:phyxminiproblems-problemphyxmini0947-vectorsireadout}
  Read a torque in newton-metre coordinates
\end{definition}
\begin{proof}
  This is the declared model definition of torque In Newton Meters.
\end{proof}
\begin{definition}[torque Magnitude In Newton Meters]
  \label{def:physics:phyx-mini-0947:phyxminiproblems-problemphyxmini0947-torquemagnitudeinnewtonmeters}
  \lean{PhyXMiniProblems.ProblemPhyXMini0947.torqueMagnitudeInNewtonMeters}
  \uses{def:physics:phyx-mini-0947:phyxminiproblems-problemphyxmini0947-torquevector, def:physics:phyx-mini-0947:phyxminiproblems-problemphyxmini0947-torqueinnewtonmeters}
  Euclidean magnitude of a torque readout, in newton-metres
\end{definition}
\begin{proof}
  This is the declared model definition of torque Magnitude In Newton Meters.
\end{proof}
\begin{definition}[spatial Cross]
  \label{def:physics:phyx-mini-0947:phyxminiproblems-problemphyxmini0947-spatialcross}
  \lean{PhyXMiniProblems.ProblemPhyXMini0947.spatialCross}
  \uses{def:physics:phyx-mini-0947:phyxminiproblems-problemphyxmini0947-spatialvector}
  Mathlib's ordinary cross product transported to EuclideanSpace ℝ (Fin 3)
\end{definition}
\begin{proof}
  This is the declared model definition of spatial Cross.
\end{proof}
\begin{definition}[x Hat]
  \label{def:physics:phyx-mini-0947:phyxminiproblems-problemphyxmini0947-xhat}
  \lean{PhyXMiniProblems.ProblemPhyXMini0947.xHat}
  \uses{def:physics:phyx-mini-0947:phyxminiproblems-problemphyxmini0947-spatialvector}
  Dimensionless positive-x coordinate unit vector
\end{definition}
\begin{proof}
  This is the declared model definition of x Hat.
\end{proof}
\begin{definition}[y Hat]
  \label{def:physics:phyx-mini-0947:phyxminiproblems-problemphyxmini0947-yhat}
  \lean{PhyXMiniProblems.ProblemPhyXMini0947.yHat}
  \uses{def:physics:phyx-mini-0947:phyxminiproblems-problemphyxmini0947-spatialvector}
  Dimensionless positive-y coordinate unit vector
\end{definition}
\begin{proof}
  This is the declared model definition of y Hat.
\end{proof}
\begin{definition}[z Hat]
  \label{def:physics:phyx-mini-0947:phyxminiproblems-problemphyxmini0947-zhat}
  \lean{PhyXMiniProblems.ProblemPhyXMini0947.zHat}
  \uses{def:physics:phyx-mini-0947:phyxminiproblems-problemphyxmini0947-spatialvector}
  Dimensionless positive-z coordinate unit vector
\end{definition}
\begin{proof}
  This is the declared model definition of z Hat.
\end{proof}
\begin{definition}[degrees To Radians]
  \label{def:physics:phyx-mini-0947:phyxminiproblems-problemphyxmini0947-degreestoradians}
  \lean{PhyXMiniProblems.ProblemPhyXMini0947.degreesToRadians}
  Convert a dimensionless degree readout to radians
\end{definition}
\begin{proof}
  This is the declared model definition of degrees To Radians.
\end{proof}
\begin{definition}[Loop Topology]
  \label{def:physics:phyx-mini-0947:phyxminiproblems-problemphyxmini0947-looptopology}
  \lean{PhyXMiniProblems.ProblemPhyXMini0947.LoopTopology}
  The one closed-loop topology stated by the problem
\end{definition}
\begin{proof}
  This is the declared model definition of Loop Topology.
\end{proof}
\begin{definition}[Coordinate Axis]
  \label{def:physics:phyx-mini-0947:phyxminiproblems-problemphyxmini0947-coordinateaxis}
  \lean{PhyXMiniProblems.ProblemPhyXMini0947.CoordinateAxis}
  Coordinate axes printed beside the loop
\end{definition}
\begin{proof}
  This is the declared model definition of Coordinate Axis.
\end{proof}
\begin{definition}[Loop Edge]
  \label{def:physics:phyx-mini-0947:phyxminiproblems-problemphyxmini0947-loopedge}
  \lean{PhyXMiniProblems.ProblemPhyXMini0947.LoopEdge}
  The four labelled geometric sides of the rectangular loop
\end{definition}
\begin{proof}
  This is the declared model definition of Loop Edge.
\end{proof}
\begin{definition}[Current Arrow Traversal]
  \label{def:physics:phyx-mini-0947:phyxminiproblems-problemphyxmini0947-currentarrowtraversal}
  \lean{PhyXMiniProblems.ProblemPhyXMini0947.CurrentArrowTraversal}
  Traversal direction of a current arrow on a particular displayed edge
\end{definition}
\begin{proof}
  This is the declared model definition of Current Arrow Traversal.
\end{proof}
\begin{definition}[Rectangular Loop Figure]
  \label{def:physics:phyx-mini-0947:phyxminiproblems-problemphyxmini0947-rectangularloopfigure}
  \lean{PhyXMiniProblems.ProblemPhyXMini0947.RectangularLoopFigure}
  \uses{def:physics:phyx-mini-0947:phyxminiproblems-problemphyxmini0947-coordinateaxis, def:physics:phyx-mini-0947:phyxminiproblems-problemphyxmini0947-loopedge, def:physics:phyx-mini-0947:phyxminiproblems-problemphyxmini0947-currentarrowtraversal}
  Literal content transcribed from the supplied 239 × 297 raster
\end{definition}
\begin{proof}
  This is the declared model definition of Rectangular Loop Figure.
\end{proof}
\begin{definition}[Rectangular Current Loop Setup]
  \label{def:physics:phyx-mini-0947:phyxminiproblems-problemphyxmini0947-rectangularcurrentloopsetup}
  \lean{PhyXMiniProblems.ProblemPhyXMini0947.RectangularCurrentLoopSetup}
  \uses{def:physics:phyx-mini-0947:phyxminiproblems-problemphyxmini0947-spatialvector, def:physics:phyx-mini-0947:phyxminiproblems-problemphyxmini0947-lengthmagnitude, def:physics:phyx-mini-0947:phyxminiproblems-problemphyxmini0947-electriccurrentmagnitude, def:physics:phyx-mini-0947:phyxminiproblems-problemphyxmini0947-magneticfluxdensitymagnitude, def:physics:phyx-mini-0947:phyxminiproblems-problemphyxmini0947-magneticfluxdensityvector, def:physics:phyx-mini-0947:phyxminiproblems-problemphyxmini0947-magneticmomentvector, def:physics:phyx-mini-0947:phyxminiproblems-problemphyxmini0947-torquevector, def:physics:phyx-mini-0947:phyxminiproblems-problemphyxmini0947-looptopology, def:physics:phyx-mini-0947:phyxminiproblems-problemphyxmini0947-coordinateaxis, def:physics:phyx-mini-0947:phyxminiproblems-problemphyxmini0947-rectangularloopfigure}
  Independent physical data for the loop and field. The magnetic moment, magnetic torque, and holding torque are independent observables constrained by governing-law structures below; none is defined from an answer choice
\end{definition}
\begin{proof}
  This is the declared model definition of Rectangular Current Loop Setup.
\end{proof}
\begin{definition}[Matches Problem Statement]
  \label{def:physics:phyx-mini-0947:phyxminiproblems-problemphyxmini0947-matchesproblemstatement}
  \lean{PhyXMiniProblems.ProblemPhyXMini0947.MatchesProblemStatement}
  \uses{def:physics:phyx-mini-0947:phyxminiproblems-problemphyxmini0947-currentinamperes, def:physics:phyx-mini-0947:phyxminiproblems-problemphyxmini0947-magneticfluxdensityinteslas, def:physics:phyx-mini-0947:phyxminiproblems-problemphyxmini0947-rectangularcurrentloopsetup}
  Numerical and qualitative data explicitly supplied in the problem prose
\end{definition}
\begin{proof}
  This is the declared model definition of Matches Problem Statement.
\end{proof}
\begin{definition}[Matches Primary Figure]
  \label{def:physics:phyx-mini-0947:phyxminiproblems-problemphyxmini0947-matchesprimaryfigure}
  \lean{PhyXMiniProblems.ProblemPhyXMini0947.MatchesPrimaryFigure}
  \uses{def:physics:phyx-mini-0947:phyxminiproblems-problemphyxmini0947-lengthinmeters, def:physics:phyx-mini-0947:phyxminiproblems-problemphyxmini0947-currentinamperes, def:physics:phyx-mini-0947:phyxminiproblems-problemphyxmini0947-degreestoradians, def:physics:phyx-mini-0947:phyxminiproblems-problemphyxmini0947-rectangularcurrentloopsetup}
  Primary-image evidence and its calibration to physical quantities. These fields contain no torque value or answer-choice conclusion
\end{definition}
\begin{proof}
  This is the declared model definition of Matches Primary Figure.
\end{proof}
\begin{definition}[Matches Rectangular Loop Geometry]
  \label{def:physics:phyx-mini-0947:phyxminiproblems-problemphyxmini0947-matchesrectangularloopgeometry}
  \lean{PhyXMiniProblems.ProblemPhyXMini0947.MatchesRectangularLoopGeometry}
  \uses{def:physics:phyx-mini-0947:phyxminiproblems-problemphyxmini0947-lengthinmeters, def:physics:phyx-mini-0947:phyxminiproblems-problemphyxmini0947-areainsquaremeters, def:physics:phyx-mini-0947:phyxminiproblems-problemphyxmini0947-spatialcross, def:physics:phyx-mini-0947:phyxminiproblems-problemphyxmini0947-xhat, def:physics:phyx-mini-0947:phyxminiproblems-problemphyxmini0947-yhat, def:physics:phyx-mini-0947:phyxminiproblems-problemphyxmini0947-zhat, def:physics:phyx-mini-0947:phyxminiproblems-problemphyxmini0947-rectangularcurrentloopsetup}
  Cartesian realization of the pictured rectangular geometry. The free-side displacement lies in the x-z plane, rotated from positive x toward negative z. Since current rises along the pivot side, the right-hand-rule area normal is yHat × freeSideDisplacementUnit. This is geometry and current orientation only; it does not mention either torque
\end{definition}
\begin{proof}
  This is the declared model definition of Matches Rectangular Loop Geometry.
\end{proof}
\begin{definition}[Models Uniform Magnetic Field]
  \label{def:physics:phyx-mini-0947:phyxminiproblems-problemphyxmini0947-modelsuniformmagneticfield}
  \lean{PhyXMiniProblems.ProblemPhyXMini0947.ModelsUniformMagneticField}
  \uses{def:physics:phyx-mini-0947:phyxminiproblems-problemphyxmini0947-magneticfluxdensityinteslas, def:physics:phyx-mini-0947:phyxminiproblems-problemphyxmini0947-magneticfluxdensityvectorinteslas, def:physics:phyx-mini-0947:phyxminiproblems-problemphyxmini0947-xhat, def:physics:phyx-mini-0947:phyxminiproblems-problemphyxmini0947-rectangularcurrentloopsetup}
  The field is spatially and temporally uniform and points along positive x. Physlib's spacetime-dependent magnetic field is calibrated by the independent dimensionful flux-density vector
\end{definition}
\begin{proof}
  This is the declared model definition of Models Uniform Magnetic Field.
\end{proof}
\begin{definition}[Satisfies Current Loop Magnetic Moment Law]
  \label{def:physics:phyx-mini-0947:phyxminiproblems-problemphyxmini0947-satisfiescurrentloopmagneticmomentlaw}
  \lean{PhyXMiniProblems.ProblemPhyXMini0947.SatisfiesCurrentLoopMagneticMomentLaw}
  \uses{def:physics:phyx-mini-0947:phyxminiproblems-problemphyxmini0947-areainsquaremeters, def:physics:phyx-mini-0947:phyxminiproblems-problemphyxmini0947-currentinamperes, def:physics:phyx-mini-0947:phyxminiproblems-problemphyxmini0947-magneticmomentinamperesquaremeters, def:physics:phyx-mini-0947:phyxminiproblems-problemphyxmini0947-rectangularcurrentloopsetup}
  The magnetic moment of a one-turn planar current loop is current times area in the current-oriented normal direction. This governing law contains no torque answer
\end{definition}
\begin{proof}
  This is the declared model definition of Satisfies Current Loop Magnetic Moment Law.
\end{proof}
\begin{definition}[Satisfies Magnetic Torque Law]
  \label{def:physics:phyx-mini-0947:phyxminiproblems-problemphyxmini0947-satisfiesmagnetictorquelaw}
  \lean{PhyXMiniProblems.ProblemPhyXMini0947.SatisfiesMagneticTorqueLaw}
  \uses{def:physics:phyx-mini-0947:phyxminiproblems-problemphyxmini0947-magneticfluxdensityvectorinteslas, def:physics:phyx-mini-0947:phyxminiproblems-problemphyxmini0947-magneticmomentinamperesquaremeters, def:physics:phyx-mini-0947:phyxminiproblems-problemphyxmini0947-torqueinnewtonmeters, def:physics:phyx-mini-0947:phyxminiproblems-problemphyxmini0947-spatialcross, def:physics:phyx-mini-0947:phyxminiproblems-problemphyxmini0947-rectangularcurrentloopsetup}
  The magnetic torque on a dipole is μ × B
\end{definition}
\begin{proof}
  This is the declared model definition of Satisfies Magnetic Torque Law.
\end{proof}
\begin{definition}[Satisfies Static Holding Condition]
  \label{def:physics:phyx-mini-0947:phyxminiproblems-problemphyxmini0947-satisfiesstaticholdingcondition}
  \lean{PhyXMiniProblems.ProblemPhyXMini0947.SatisfiesStaticHoldingCondition}
  \uses{def:physics:phyx-mini-0947:phyxminiproblems-problemphyxmini0947-torqueinnewtonmeters, def:physics:phyx-mini-0947:phyxminiproblems-problemphyxmini0947-rectangularcurrentloopsetup}
  To hold the loop fixed, the applied holding torque balances the magnetic torque. This is the static-equilibrium law, not the requested evaluated holding torque
\end{definition}
\begin{proof}
  This is the declared model definition of Satisfies Static Holding Condition.
\end{proof}
\begin{definition}[Points Along Positive Axis]
  \label{def:physics:phyx-mini-0947:phyxminiproblems-problemphyxmini0947-pointsalongpositiveaxis}
  \lean{PhyXMiniProblems.ProblemPhyXMini0947.PointsAlongPositiveAxis}
  \uses{def:physics:phyx-mini-0947:phyxminiproblems-problemphyxmini0947-spatialvector, def:physics:phyx-mini-0947:phyxminiproblems-problemphyxmini0947-xhat, def:physics:phyx-mini-0947:phyxminiproblems-problemphyxmini0947-yhat, def:physics:phyx-mini-0947:phyxminiproblems-problemphyxmini0947-zhat, def:physics:phyx-mini-0947:phyxminiproblems-problemphyxmini0947-coordinateaxis}
  A vector points strictly along the positive direction of an axis
\end{definition}
\begin{proof}
  This is the declared model definition of Points Along Positive Axis.
\end{proof}
\begin{definition}[Answer Choice]
  \label{def:physics:phyx-mini-0947:phyxminiproblems-problemphyxmini0947-answerchoice}
  \lean{PhyXMiniProblems.ProblemPhyXMini0947.AnswerChoice}
  Labels of the four displayed torque-magnitude choices
\end{definition}
\begin{proof}
  This is the declared model definition of Answer Choice.
\end{proof}
\begin{definition}[displayed Torque Magnitude In Newton Meters]
  \label{def:physics:phyx-mini-0947:phyxminiproblems-problemphyxmini0947-displayedtorquemagnitudeinnewtonmeters}
  \lean{PhyXMiniProblems.ProblemPhyXMini0947.displayedTorqueMagnitudeInNewtonMeters}
  \uses{def:physics:phyx-mini-0947:phyxminiproblems-problemphyxmini0947-answerchoice}
  Torque magnitude printed beside each choice, in newton-metres
\end{definition}
\begin{proof}
  This is the declared model definition of displayed Torque Magnitude In Newton Meters.
\end{proof}
\begin{definition}[recorded Answer Choice]
  \label{def:physics:phyx-mini-0947:phyxminiproblems-problemphyxmini0947-recordedanswerchoice}
  \lean{PhyXMiniProblems.ProblemPhyXMini0947.recordedAnswerChoice}
  \uses{def:physics:phyx-mini-0947:phyxminiproblems-problemphyxmini0947-answerchoice}
  The answer label recorded in the dataset metadata
\end{definition}
\begin{proof}
  This is the declared model definition of recorded Answer Choice.
\end{proof}
\begin{definition}[Agrees At Millinewton Meter Precision]
  \label{def:physics:phyx-mini-0947:phyxminiproblems-problemphyxmini0947-agreesatmillinewtonmeterprecision}
  \lean{PhyXMiniProblems.ProblemPhyXMini0947.AgreesAtMillinewtonMeterPrecision}
  \uses{def:physics:phyx-mini-0947:phyxminiproblems-problemphyxmini0947-torquemagnitudeinnewtonmeters, def:physics:phyx-mini-0947:phyxminiproblems-problemphyxmini0947-rectangularcurrentloopsetup, def:physics:phyx-mini-0947:phyxminiproblems-problemphyxmini0947-answerchoice, def:physics:phyx-mini-0947:phyxminiproblems-problemphyxmini0947-displayedtorquemagnitudeinnewtonmeters}
  Agreement after rounding a torque magnitude to the nearest 0.001 N m
\end{definition}
\begin{proof}
  This is the declared model definition of Agrees At Millinewton Meter Precision.
\end{proof}
% --- Archon declaration topology end ---
% --- Archon physics formalization source end ---
